% figures/fig2/fig2.tex — auto-generated by src/adaptive_behavior/make_fig_tex.py
%
% Layout driven by exact PowerPoint coordinates from layout.tex.
% To regenerate this file after editing the pptx:
%   python -m adaptive_behavior.update_layout fig2
%   python -m adaptive_behavior.make_fig_tex fig2 --overwrite

% figures/fig2/layout.tex � auto-generated, do not edit by hand
% Regenerate: python -m adaptive_behavior.extract_pptx_layout figures/fig2/fig2.pptx --figure fig2

% Per-panel positions and sizes (exact PowerPoint measurements in mm)
\newlength{\FigIIAX}\setlength{\FigIIAX}{0.0mm}
\newlength{\FigIIAY}\setlength{\FigIIAY}{0.0mm}
\newlength{\FigIIAW}\setlength{\FigIIAW}{25.7mm}
\newlength{\FigIIAH}\setlength{\FigIIAH}{25.7mm}
\newlength{\FigIIBX}\setlength{\FigIIBX}{27.2mm}
\newlength{\FigIIBY}\setlength{\FigIIBY}{0.3mm}
\newlength{\FigIIBW}\setlength{\FigIIBW}{52.9mm}
\newlength{\FigIIBH}\setlength{\FigIIBH}{25.2mm}
\newlength{\FigIICX}\setlength{\FigIICX}{0.6mm}
\newlength{\FigIICY}\setlength{\FigIICY}{26.9mm}
\newlength{\FigIICW}\setlength{\FigIICW}{25.7mm}
\newlength{\FigIICH}\setlength{\FigIICH}{25.4mm}
\newlength{\FigIIDX}\setlength{\FigIIDX}{27.2mm}
\newlength{\FigIIDY}\setlength{\FigIIDY}{27.1mm}
\newlength{\FigIIDW}\setlength{\FigIIDW}{12.3mm}
\newlength{\FigIIDH}\setlength{\FigIIDH}{25.3mm}
\newlength{\FigIIEX}\setlength{\FigIIEX}{40.6mm}
\newlength{\FigIIEY}\setlength{\FigIIEY}{27.4mm}
\newlength{\FigIIEW}\setlength{\FigIIEW}{25.0mm}
\newlength{\FigIIEH}\setlength{\FigIIEH}{25.0mm}
\newlength{\FigIIFX}\setlength{\FigIIFX}{67.0mm}
\newlength{\FigIIFY}\setlength{\FigIIFY}{27.5mm}
\newlength{\FigIIFW}\setlength{\FigIIFW}{12.8mm}
\newlength{\FigIIFH}\setlength{\FigIIFH}{24.9mm}
\newlength{\FigIIGX}\setlength{\FigIIGX}{0.8mm}
\newlength{\FigIIGY}\setlength{\FigIIGY}{53.4mm}
\newlength{\FigIIGW}\setlength{\FigIIGW}{20.6mm}
\newlength{\FigIIGH}\setlength{\FigIIGH}{25.7mm}
\newlength{\FigIIHX}\setlength{\FigIIHX}{22.7mm}
\newlength{\FigIIHY}\setlength{\FigIIHY}{53.5mm}
\newlength{\FigIIHW}\setlength{\FigIIHW}{20.6mm}
\newlength{\FigIIHH}\setlength{\FigIIHH}{25.4mm}
\newlength{\FigIIIX}\setlength{\FigIIIX}{44.2mm}
\newlength{\FigIIIY}\setlength{\FigIIIY}{53.9mm}
\newlength{\FigIIIW}\setlength{\FigIIIW}{21.4mm}
\newlength{\FigIIIH}\setlength{\FigIIIH}{25.0mm}
\newlength{\FigIIJX}\setlength{\FigIIJX}{66.4mm}
\newlength{\FigIIJY}\setlength{\FigIIJY}{54.0mm}
\newlength{\FigIIJW}\setlength{\FigIIJW}{18.9mm}
\newlength{\FigIIJH}\setlength{\FigIIJH}{24.9mm}
\newlength{\FigIIKX}\setlength{\FigIIKX}{1.2mm}
\newlength{\FigIIKY}\setlength{\FigIIKY}{80.1mm}
\newlength{\FigIIKW}\setlength{\FigIIKW}{25.1mm}
\newlength{\FigIIKH}\setlength{\FigIIKH}{25.7mm}
\newlength{\FigIILX}\setlength{\FigIILX}{27.2mm}
\newlength{\FigIILY}\setlength{\FigIILY}{80.3mm}
\newlength{\FigIILW}\setlength{\FigIILW}{27.9mm}
\newlength{\FigIILH}\setlength{\FigIILH}{25.7mm}
\newlength{\FigIIMX}\setlength{\FigIIMX}{56.6mm}
\newlength{\FigIIMY}\setlength{\FigIIMY}{80.3mm}
\newlength{\FigIIMW}\setlength{\FigIIMW}{28.7mm}
\newlength{\FigIIMH}\setlength{\FigIIMH}{25.2mm}
\newlength{\FigIINX}\setlength{\FigIINX}{1.2mm}
\newlength{\FigIINY}\setlength{\FigIINY}{106.8mm}
\newlength{\FigIINW}\setlength{\FigIINW}{25.1mm}
\newlength{\FigIINH}\setlength{\FigIINH}{24.8mm}
\newlength{\FigIIOX}\setlength{\FigIIOX}{27.3mm}
\newlength{\FigIIOY}\setlength{\FigIIOY}{106.7mm}
\newlength{\FigIIOW}\setlength{\FigIIOW}{27.8mm}
\newlength{\FigIIOH}\setlength{\FigIIOH}{24.9mm}
\newlength{\FigIIPX}\setlength{\FigIIPX}{56.2mm}
\newlength{\FigIIPY}\setlength{\FigIIPY}{107.1mm}
\newlength{\FigIIPW}\setlength{\FigIIPW}{29.1mm}
\newlength{\FigIIPH}\setlength{\FigIIPH}{24.0mm}

% Derived group widths (bounding boxes of panel clusters)
\newlength{\FigIIRightW}\setlength{\FigIIRightW}{28.7mm}
\newlength{\FigIILeftSubW}\setlength{\FigIILeftSubW}{-0.9mm}
\newlength{\FigIIRightSubW}\setlength{\FigIIRightSubW}{-45.6mm}


\begin{figure}[H]
  \vspace{-13mm}   % pulls the figure up; adjust to taste
  \centering  % centered
  \captionsetup{font=footnotesize}  % 10 pt caption (footnotesize at 12 pt base)
  \setcounter{subfigure}{0}
  \makebox[\linewidth][c]{\makebox[195mm][c]{%  % natural size (no scaling), centered
  \begin{tikzpicture}[inner sep=0pt, outer sep=0pt]

    \node[anchor=north west] at (\FigIIAX, -\FigIIAY) {%
      \begin{subfigure}[t]{\FigIIAW}%
        \panelgraphic[width=\linewidth]{figures/fig2/panels/fig2_A.pdf}%
        \phantomsubcaption\label{fig:2-a}%
        \phantomsubcaption\label{fig:2-b}%
        \phantomsubcaption\label{fig:2-c}%
        \phantomsubcaption\label{fig:2-d}%
        \phantomsubcaption\label{fig:2-e}%
        \phantomsubcaption\label{fig:2-f}%
        \phantomsubcaption\label{fig:2-g}%
        \phantomsubcaption\label{fig:2-h}%
        \phantomsubcaption\label{fig:2-i}%
        \phantomsubcaption\label{fig:2-j}%
        \phantomsubcaption\label{fig:2-k}%
        \phantomsubcaption\label{fig:2-l}%
        \phantomsubcaption\label{fig:2-m}%
        \phantomsubcaption\label{fig:2-n}%
        \phantomsubcaption\label{fig:2-o}%
        \phantomsubcaption\label{fig:2-p}%
        \phantomsubcaption\label{fig:2-q}%
        \phantomsubcaption\label{fig:2-r}%
      \end{subfigure}};

  \end{tikzpicture}}}% end makebox
\end{figure}

\addtocounter{figure}{-1}  % the image-only figure above steps the counter

\begin{figure}[H]
  \captionsetup{font=footnotesize}  % 10 pt caption (footnotesize at 12 pt base)
  \caption{\textbf{Patch-leaving behavior is consistent with a Bayesian formulation of the Marginal Value Theorem.}
    (\subref{fig:2-a}) MVT predicts that optimal patch leaving is set by the global reward rate (slope of black tangent lines), and foragers should leave high or low value odor patch types (orange and green lines) at the same local reward rate. 
    (\subref{fig:2-b}) Reward probability at leaving is matched across patch types, while rewards collected, and number of reward omissions experienced differed significantly. Lines are individual mice ($n=24$).
    (\subref{fig:2-c}) Mice have odor-cued expectations about patch quatlity. Given that the first odor site in a patch was unrewarded, mice were more likely to leave the low value patch (green) than the high. 
    (\subref{fig:2-d}) Change in $P(\text{reward})$ at leaving relative to control as a function of inter-patch travel distance. Longer distances led mice to leave at lower reward probabilities, while shorter distances produced the opposite effect, consistent with MVT predictions.
    (\subref{fig:2-e},\subref{fig:2-f} Experiments manipulating global reward rate.  In (\subref{fig:2-e}), odor 2 reward statistics were increased to match those of odor 1; which led to an increase in global reward. Relative to control sessions, both patch types shifted to higher reward probabilities (left), meaning mice left patches earlier. To achieve this, rewards collected shifted in opposite directions across patch types (right) .In (\subref{fig:2-f}) Opposite manipulating where we decreased the global reward rate. Both patch types shifted to lower reward probabilities meaning the mice left patches later. Both results are qualitatively compatible with MVT. 
    (\subref{fig:2-g}) Bayesian MVT model schematic. The agent maintains a separate posterior over reward probability $q$ for each odor-cued patch type, updated after each binary outcome via Bayes' rule with a forgetting parameter $\delta < 1$ that discounts accumulated prior beliefs at each update step. At each odor site, a sample $\tilde{q}$ is drawn from the current posterior; the agent leaves if the implied local reward rate falls below a running estimate of the global reward rate $\bar{R}$.
    (\subref{fig:2-h}--\subref{fig:2-l}) Bayesian MVT model simulations reproducing the mouse results in (\subref{fig:2-b}--\subref{fig:2-f}): reward probability at leaving, rewards collected and consecutive reward omissions (\subref{fig:2-h}), probability of leaving after early consecutive omissions (\subref{fig:2-i}, change in $P(\text{reward})$ at leaving with inter-patch distance (\subref{fig:2-j}), and results for higher (\subref{fig:2-q}) and lower (\subref{fig:2-l}) global reward rate environments. 
    Significance (paired $t$-tests across subjects): ns, not significant; $^{*}p<0.05$, $^{**}p<0.01$, $^{***}p<0.001$.}
  \label{fig:2}
\end{figure}
